\documentclass[10pt, a4paper]{article}
\usepackage{preamble}

\title{Algebra II \\
    \large Cheat Sheet}
\author{Luke Phillips}
\date{October 2025}

\begin{document}

\maketitle

\newpage

\section{Cheat Sheet}

\begin{definition}
    A ring $R$ is a set with two binary operations $+, \cdot$ such that $+$ makes $R$ into an abelian group and such that
    \begin{enumerate}[label = (\roman*)]
        \item (Identity)
        \[
        \exists 1 \in R, 1 \cdot x = x \cdot 1 = x,\quad\forall x \in R.
        \]

        \item (Associativity)
        \[
        x(yz) = (xy)z,\quad\forall x, y, z \in R.
        \]

        \item (Distributivity)
        \[
        x(y + z) = xy + xz\text{ and } (y + z)x = yx + zx,\quad\forall x, y, z \in R.
        \]
    \end{enumerate}
\end{definition}

The \textbf{modulo notation} for the set of integers modulo $n$ is denoted $\Z / n$ called the residue class.

\begin{definition}
    A subring $S$ of a ring $(R, +, \cdot)$ is a set $S \subseteq R$ such that
    \begin{enumerate}[label = (\roman*)]
        \item $0 \in S, 1 \in S$.

        \item $a, b \in S \implies a + b \in S$.

        \item $a, b \in S \implies ab \in S$.

        \item $a \in S \implies -a \in S$.
    \end{enumerate}
\end{definition}
\textit{Note:
a subring is something that has $0$ and $1$,
and that you can add or multiply any elements in it together to get something still in it
(as well as negation).}

\begin{definition}
    A ring $(R, +, \cdot)$ is a field if
    \begin{enumerate}[label = (\roman*)]
        \item $R$ is commutative.

        \item $1 \neq 0$
        ($R$ has at least two elements).

        \item For every $a \in R$,
        such that $a \neq 0, \exists b \in R$ such that $ab = 1$
        (i.e. for every $a$ there is an inverse $b = a ^ {-1}$).
    \end{enumerate}
\end{definition}

\begin{remark}
    Inverses are \textbf{unique},
    so if you want to prove something about an inverse this could be a good place to start.
\end{remark}

\vspace{1em}

\begin{lemma}[Well-Ordering Principle
    (WOP)]
    Let $S$ be a non-empty subset of $\N_0 = \N \cup \{0\}$.
    Then $S$ has a smallest element.
\end{lemma}

\vspace{0.5em}

\begin{remark}
    The \textbf{above} could be used in a proof if we need the \textbf{smallest} $m$ in a finite set that satisfies a condition.
\end{remark}

\vspace{1em}

\begin{proposition}
    Let $a, b \in \Z$,
    $b \neq 0$.
    Then there exist $q, r \in \Z$ such that 
    \[
    a = q \cdot b + r\qquad\text{and}\qquad 0 \leq r < |b|.
    \]
\end{proposition}

\begin{remark}
    This means that for $a \in \Z$ we can always split it up into something with a multiple that we choose and a remainder,
    quite useful in \textbf{modulo arithmetic}.
\end{remark}

\begin{theorem}
    Let $a, b \in \Z$,
    not both $0$.
    Then $\gcd(a, b)$ exists and can be calculated by the Euclidean algorithm.
    Moreover,
    if $d = \gcd(a, b)$ then $\exists x, y \in \Z$ such that
    \[
    ax + by = d.
    \]
\end{theorem}

\begin{lemma}
    If $p$ prime and $p \mid ab$,
    $a, b \in \Z$,
    $p \mid a$ or $p \mid b$.
\end{lemma}

\begin{definition}
    Let $R$ be a ring.
    An element $a \in R$ is a unit
    (invertible)
    if $\exists b \in R$ such that $ab = ba = 1$.
\end{definition}

The \textbf{set of units} of $R$ is denoted $R ^ {\times}$.

\begin{definition}
    Let $R$ be a ring.
    $r \in R$ is irreducible if it satisfies the following:
    \begin{enumerate}[label = (\roman*)]
        \item $r$ is not a unit.

        \item If $r = ab$,
        either $a$ or $b$ is a unit.
    \end{enumerate}
\end{definition}

\begin{lemma}
    Let $R$
    (a ring)
    be $\Z$ or $F[x]$,
    $F$ a field.
    If $p \in R$ is irreducible and $p \mid ab$,
    $a, b \in \R$,
    then $p \mid a$ or $p \mid b$.
\end{lemma}

\begin{proposition}
    An element $\bar{a} \in \Z / n$ is a unit if and only if $\gcd(a, n) = 1$.
\end{proposition}

\begin{corollary}
    $\Z / p$ is a field if and only if $p$ is prime.
\end{corollary}

\begin{definition}
    Let $R$ be a ring.
    If $a \in R$ is such that $\exists b \neq 0, ab = 0$ then $a$ is a zero divisor.
\end{definition}

\begin{definition}
    Let $R$ be a commutative ring with $0 \neq 1$,
    and no zero divisors apart from $0$,
    it is called an integral domain.
\end{definition}

The \textbf{essential property} of integral domains is cancellation:
\[
ax = bx,\qquad x \neq 0 \implies a = b.
\]

\begin{lemma}
    $\Z / n$ is an integral domain if and only if $n$ is a prime
\end{lemma}

\vspace{0.5em}

\begin{remark}
    From the above
    (and a previous corollary)
    we can see that $\Z / p$ is an integral domain and a field if and only if $p$ is a prime.
\end{remark}

\vspace{1em}

\begin{proposition}
    If $R$ is an integral domain,
    then $R[x]$ is an integral domain.
\end{proposition}

\begin{definition}
    Let $R$ be a commutative ring.
    $x \in R$ is called prime if:
    \begin{enumerate}[label = (\roman*)]
        \item $x \neq 0$ and $x \notin R ^ {\times}$.

        \item If $x \mid ab$,
        $a, b \in R$,
        then $x \mid a$ or $x \mid b$.
    \end{enumerate}
\end{definition}

\begin{lemma}
    Let $R$ be an integral domain and $x \in R$ a prime.
    Then $x$ is irreducible.
\end{lemma}

\begin{definition}
    An integral domain $R$ is called a unique factorisation domain
    (UFD)
    if every non-zero non-unit can be written as a product of irreducible elements and this product is unique up to order of the factor and units.
\end{definition}

\begin{proposition}
    Let $F$ be a field and $f(x) \in F[x]$.

    \begin{enumerate}[label = (\roman*)]
        \item If $\deg{f} = 1$,
        $f(x)$ is irreducible.

        \item If $\deg{f} = 2$ or $3$,
        $f(x)$ is irreducible if and only if $f(x)$ has no roots in $F$.

        \item If $\deg{f} = 4$ then $f(x)$ is irreducible if and only if it has no roots in $F$ and is not the product of two quadratic polynomials.
    \end{enumerate}
\end{proposition}

\begin{proposition}[Rational root test]
    
\end{proposition}


\end{document}